\label{sec:bias}

\subsection{Jefferson Maximises Large-State Representation}

\begin{theorem}[Jefferson Large-State Bias]
  \label{thm:jefferson-bias}
  Among all divisor methods, Jefferson is the most large-state-favoring:
  for any partition of states into ``large'' ($p_i > p_{\text{med}}$)
  and ``small'' ($p_i \leq p_{\text{med}}$), Jefferson maximises the
  aggregate seat share of large states.
\end{theorem}

\begin{proof}[Proof sketch]
  For floor rounding with divisor $d$, each state receives $\lfloor p_i/d \rfloor$
  seats.
  The fraction of seats ``lost'' by state $i$ to rounding is
  $(p_i/d - \lfloor p_i/d \rfloor)/(\lfloor p_i/d \rfloor + (p_i/d - \lfloor p_i/d \rfloor))$.
  For small $s_i$ (small states), this fraction is large ---
  the fractional part is a large share of the seat count.
  For large $s_i$ (large states), the fractional part is a small
  share of the seat count.
  Jefferson discards all fractional parts; any divisor method that
  preserves some fractional parts (HH, Webster, Adams) redistributes
  those fractional parts from large states to small states,
  reducing large states' share.
  Hence Jefferson, which discards all fractional parts equally,
  effectively transfers the most seat-share to large states
  relative to any other divisor method.
  See \citet{balinski1982fair}, ch.~4, for the formal ordering proof.
\end{proof}

\subsection{Quantitative Bias Characterisation}

For each divisor method $M$, define the \emph{relative deviation}
from proportionality for state $i$ as:
\[
  \delta_i^M = \frac{s_i^M/H}{p_i/N} - 1,
\]
where $s_i^M$ is the seat count under method $M$, $H = 435$,
and $N$ is total population.
$\delta_i^M > 0$ means overrepresentation; $\delta_i^M < 0$ means
underrepresentation.

Under Jefferson for 2020 Census data:
\begin{itemize}
  \item \textbf{Wyoming}: exact quota $\approx 0.758$; receives 1 seat
    (constitutional floor).
    $\delta_{\text{WY}}^{\text{Jeff}} = (1/435) / (576{,}851/331{,}108{,}434) - 1
    \approx +32\%$ (massively overrepresented due to constitutional floor,
    not Jefferson's rounding).
  \item \textbf{California}: exact quota $\approx 51.95$; receives 51 seats
    under Jefferson.
    $\delta_{\text{CA}}^{\text{Jeff}} = (51/435) / (39{,}538{,}223/331{,}108{,}434) - 1
    \approx -2\%$ (slightly underrepresented by rounding).
  \item \textbf{Texas}: exact quota $\approx 38.31$; receives 38 seats.
    $\delta_{\text{TX}}^{\text{Jeff}} \approx -0.8\%$.
\end{itemize}
Large states (CA, TX, FL) are slightly underrepresented by floor rounding
because their large fractional parts (0.95, 0.31, etc.) are discarded.
Small states receive exactly 1 seat (constitutional floor),
regardless of their fractional quotas below 1.

The large-state-favoring bias of Jefferson is most visible not in the
constitutional-floor states (where all methods agree) but in
mid-size states near a rounding boundary.
A state with exact quota 3.9 would receive 3 seats under Jefferson
but 4 seats under Adams (and 4 under HH and Webster if quota $> 3.5$).
The loss of 0.9 fractional seats is a 23\% reduction in seat count.

\subsection{Comparison to Huntington-Hill for 2020}

For the 2020 Census, Jefferson and Huntington-Hill differ for
three states:
\begin{itemize}
  \item States that \emph{lose} seats under Jefferson relative to HH:
    Minnesota (quota $\approx 7.50$: HH gives 8, Jefferson gives 7
    because floor of 7.50 = 7).
  \item States that \emph{gain} seats under Jefferson relative to HH:
    Texas (quota $\approx 38.31$: both give 38).
    California (quota $\approx 51.95$: Jefferson gives 51, HH gives 52).
\end{itemize}
(These figures are approximate; the exact Jefferson divisor must be
found by binary search and differs from the HH divisor.)
The large-state-favoring bias of Jefferson manifests as:
California and Texas lose 1 seat each under Jefferson relative to HH
(because HH's lower threshold for rounding up awards their fractional
excess seats), while Minnesota gains 1 seat under HH (because HH's
threshold is below 0.5, capturing Minnesota's near-0.5 fractional part).
